%!TEX root = ../main.tex
\section{Proof of Theorem~\ref{thm:lqg-escape-ball} }
\label{sec:app-proof-escape-ball}

\begin{proof}
Fix \(R>0\). Define
\[
m(\sigma):=\sigma-\eta g(\sigma).
\]
Since \(g\) is continuous at \(0\) and \(g(0)=0\), we have
\[
m(\sigma)\to 0
\qquad \text{as } \sigma\to 0.
\]
Hence, for all sufficiently small \(\delta>0\),
\[
B_\delta:=\sup_{|\sigma|\le \delta}|m(\sigma)|<\infty,
\]
and moreover
\[
B_\delta\to 0
\qquad \text{as } \delta\downarrow 0.
\]

Choose \(\delta>0\) sufficiently small so that
\[
\frac{c}{\delta^2}>
\left(\frac{R+B_\delta}{\eta}\right)^2.
\]
This is possible because \(c/\delta^2\to\infty\) as \(\delta\downarrow0\), while
\(B_\delta\to0\).

Now fix \(k\) and suppose that
\[
0<|\sigma_k|\le \delta.
\]
By the assumed lower bound on the conditional second moment,
\[
\mathbb E[\varepsilon_{k+1}^2\mid \sigma_k]
\ge
\frac{c}{\sigma_k^2}
\ge
\frac{c}{\delta^2}
>
\left(\frac{R+B_\delta}{\eta}\right)^2.
\]
Therefore, it cannot be true that
\[
|\varepsilon_{k+1}|
\le
\frac{R+B_\delta}{\eta}
\qquad \text{with conditional probability }1
\text{ given } \sigma_k,
\]
because otherwise we would have
\[
\mathbb E[\varepsilon_{k+1}^2\mid \sigma_k]
\le
\left(\frac{R+B_\delta}{\eta}\right)^2,
\]
contradicting the preceding strict inequality. Hence
\[
\mathbb P\!\left(
|\varepsilon_{k+1}|>\frac{R+B_\delta}{\eta}
\,\middle|\,
\sigma_k
\right)>0.
\]

On the event
\[
\left\{
|\varepsilon_{k+1}|>\frac{R+B_\delta}{\eta}
\right\},
\]
we have, using the update rule,
\[
|\sigma_{k+1}|
=
|m(\sigma_k)-\eta\varepsilon_{k+1}|
\ge
\eta|\varepsilon_{k+1}|-|m(\sigma_k)|.
\]
Since \(|\sigma_k|\le \delta\), we have \(|m(\sigma_k)|\le B_\delta\). Therefore,
\[
|\sigma_{k+1}|
>
(R+B_\delta)-B_\delta
=
R.
\]
Thus,
\[
\mathbb P\bigl(|\sigma_{k+1}|>R\mid \sigma_k\bigr)>0.
\]
This proves the claim.
\end{proof}